\section{Embedding into Larger Dimensions and Proof of Conjecture~3}

\subsection*{Part (a): Block-diagonal embedding preserves $\rho_r$}

\begin{lemma}\label{lem:block-norm}
For every $A\in\mathbb{R}^{m\times m}$, every integer $s \geq 0$, and every real $r \geq 1$:
\[
\|\operatorname{diag}(A, 0_s)\|_{r\to r} = \|A\|_{r\to r}.
\]
\end{lemma}

\begin{proof}
Write $x = (x_1, x_2) \in \mathbb{R}^m \oplus \mathbb{R}^s$. Then $\|x\|_r^r = \|x_1\|_r^r + \|x_2\|_r^r$, so $\|x_1\|_r \leq \|x\|_r$.

\textit{Upper bound:} $\|\operatorname{diag}(A,0_s)x\|_r = \|(Ax_1, 0)\|_r = \|Ax_1\|_r \leq \|A\|_{r\to r}\|x_1\|_r \leq \|A\|_{r\to r}\|x\|_r$. Taking the supremum over $x \neq 0$ gives $\|\operatorname{diag}(A,0_s)\|_{r\to r} \leq \|A\|_{r\to r}$.

\textit{Lower bound:} Since the $\ell^r$-unit sphere in $\mathbb{R}^m$ is compact and $u \mapsto \|Au\|_r$ is continuous, the supremum $\|A\|_{r\to r} = \sup_{\|u\|_r=1}\|Au\|_r$ is attained at some $u^* \in \mathbb{R}^m$. Then
\[
\|\operatorname{diag}(A,0_s)\|_{r\to r} \geq \frac{\|(Au^*, 0)\|_r}{\|(u^*, 0)\|_r} = \frac{\|Au^*\|_r}{\|u^*\|_r} = \|A\|_{r\to r}. \qedhere
\]
\end{proof}

\begin{proposition}
For every integer $n \geq m$ and every real $r \geq 1$: $\rho_r(M_n) = \rho_r(H)$, where $M_n = \operatorname{diag}(H, 0_{n-m})$.
\end{proposition}

\begin{proof}
Any positive diagonal $D \in \mathbb{R}^{n\times n}$ can be written as $D = \operatorname{diag}(D_1, D_2)$ with $D_1$ and $D_2$ positive diagonal. Then $DM_nD^{-1} = \operatorname{diag}(D_1HD_1^{-1}, 0_{n-m})$. By Lemma~\ref{lem:block-norm}, $\|DM_nD^{-1}\|_{r\to r} = \|D_1HD_1^{-1}\|_{r\to r}$.

As $D$ ranges over all positive diagonal $n\times n$ matrices, $D_1$ ranges over all positive diagonal $m\times m$ matrices (independently of $D_2$). Hence
\[
\rho_r(M_n) = \inf_D \|DM_nD^{-1}\|_{r\to r} = \inf_{D_1} \|D_1HD_1^{-1}\|_{r\to r} = \rho_r(H). \qedhere
\]
\end{proof}

\subsection*{Part (b): Proof of Conjecture~3}

We use the Sylvester--Hadamard matrices $H_4$ and $H_8$ constructed via the recursion in Section~1. Explicitly,
\[
H_4 = \begin{pmatrix} 1&1&1&1\\ 1&-1&1&-1\\ 1&1&-1&-1\\ 1&-1&-1&1 \end{pmatrix}.
\]
We verify $H_4 = H_4^\top$ and $H_4^2 = 4I_4$ directly. First, $H_4$ is visibly symmetric. For $H_4^2$: the $(i,i)$-entry of $H_4^2$ equals $\sum_j H_{ij}^2 = 4$ (since each entry is $\pm 1$); and for $i\neq k$, the $(i,k)$-entry equals the inner product of rows $i$ and $k$:
\begin{align*}
\text{rows 1,2:}&\quad 1\cdot1+1\cdot(-1)+1\cdot1+1\cdot(-1)=0, \\
\text{rows 1,3:}&\quad 1\cdot1+1\cdot1+1\cdot(-1)+1\cdot(-1)=0, \\
\text{rows 1,4:}&\quad 1\cdot1+1\cdot(-1)+1\cdot(-1)+1\cdot1=0, \\
\text{rows 2,3:}&\quad 1\cdot1+(-1)\cdot1+1\cdot(-1)+(-1)\cdot(-1)=0, \\
\text{rows 2,4:}&\quad 1\cdot1+(-1)\cdot(-1)+1\cdot(-1)+(-1)\cdot1=0, \\
\text{rows 3,4:}&\quad 1\cdot1+1\cdot(-1)+(-1)\cdot(-1)+(-1)\cdot1=0.
\end{align*}
Hence $H_4^2 = 4I_4$, and since $H_4 = H_4^\top$, also $H_4^\top H_4 = 4I_4$. Thus $H_4$ is the Sylvester--Hadamard matrix of order $4 = 2^2$, and Theorem~\ref{thm:rho-formula} applies to give $\rho_q(H_4) = 4^{1-1/q}$ for every $q \geq 2$.

Next, define $H_8$ via the Sylvester construction:
\[
H_8 = \begin{pmatrix} H_4 & H_4 \\ H_4 & -H_4 \end{pmatrix}.
\]
By Section~1 (with $k=3$, $m=8$, $n=4$), $H_8$ is the Sylvester--Hadamard matrix of order $8=2^3$, satisfying $H_8 = H_8^\top$ and $H_8^2 = 8I_8$. Theorem~\ref{thm:rho-formula} gives $\rho_q(H_8) = 8^{1-1/q}$ for every $q \geq 2$.

\medskip

Let $n > 6$ and $p > 2$ be integers ($p \geq 3$).

\textbf{Case 1: $n = 7$.} Take $M = \operatorname{diag}(H_4, 0_3) \in \mathbb{R}^{7\times 7}$. By Part~(a), $\rho_p(M) = \rho_p(H_4) = 4^{1-1/p}$ and $\rho_{p+1}(M) = 4^{1-1/(p+1)}$. Since $1-1/p < 1-1/(p+1)$ and $4 > 1$:
\[
\rho_p(M) = 4^{1-1/p} < 4^{1-1/(p+1)} = \rho_{p+1}(M).
\]

\textbf{Case 2: $n \geq 8$.} Take $M = \operatorname{diag}(H_8, 0_{n-8}) \in \mathbb{R}^{n\times n}$ (with the zero block omitted when $n = 8$). By Part~(a), $\rho_p(M) = 8^{1-1/p}$ and $\rho_{p+1}(M) = 8^{1-1/(p+1)}$. Since $8 > 1$:
\[
\rho_p(M) = 8^{1-1/p} < 8^{1-1/(p+1)} = \rho_{p+1}(M).
\]

\medskip

In both cases, for every integer $n > 6$ and every integer $p > 2$, we have exhibited $M \in \mathbb{R}^{n\times n}$ with $\rho_p(M) < \rho_{p+1}(M)$. This proves Conjecture~3. \qed
